\begin{proof}[Proof of \cref{obj:thm:unrestricted-large-alphabet-deficit-lower}]
\begin{enumerate}
\item Put
\[
\delta:=1-q\in[0,1).
\]
Since \(d\ge1024n^2\ge1\), the hypotheses of
\cref{obj:thm:unrestricted-near-complete-arrival-envelope} hold. Its lower bound and the last term of its upper envelope give
\[
\frac1{256n}\le\mathfrak R^{+}_{n,d,q},
\qquad
\mathfrak R^{+}_{n,d,q}\le\frac4n+\delta^2.
\]
If \(\delta^2\le n^{-1}\), then
\[
\frac{\delta^2}{1024}
\le\frac1{1024n}
\le\frac1{256n}
\le\mathfrak R^{+}_{n,d,q}.
\]
This case includes \(q=1\). For the remaining case, assume
\[
\delta^2>n^{-1},
\qquad\text{so in particular}\qquad n^{-1}\le\delta^2.
\]
We establish the deficit lower bound in this regime. % lean: CausalSmith.Stat.MarRareqLogfrontier.unrestricted_large_alphabet_deficit_lower

\item Write an ordered observed sample as
\[
\mathbf o=(o_1,\ldots,o_n)
\in\left(\{0,\ldots,d-1\}\times\{0,1\}^4\right)^n,
\]
with each record ordered as \((X,A,S,R,RY)\). The decision class contains the constant zero kernel
\[
\mathsf T_0(\mathbf o,\cdot):=\operatorname{Dirac}(0).
\]
Fix an arbitrary \(\mathsf T\in\mathcal T_n\), the bounded, parameter-independent kernel class of
\cref{obj:def:unrestricted-minimax-risk}, and define its risk under a full-data law \(P\) by
\[
\mathcal R_{\mathsf T}(P)
:=
\int\!\int_{[-1,1]}
(w-\tau(P))^2\,
\mathsf T(\mathbf o,dw)\,
\mathcal L_P(\mathbf O_n)(d\mathbf o).
\]
Define the separation between the two target centers by
\[
h_{\mathrm{sep}}
:=\frac12-\frac q{1+q}
=\frac{\delta}{2(1+q)}.
\]
Because \(1<1+q\le2\) and \(\delta\ge0\), this identity gives
\[
\frac{\delta}{4}\le h_{\mathrm{sep}}\le\frac{\delta}{2},
\qquad h_{\mathrm{sep}}\ge0.
\]

Construct a comparison law \(P_\star\) by taking a uniform baseline, zero surrogates and control outcome, and a Bernoulli treated potential outcome:
\[
X\sim\operatorname{Unif}\{0,\ldots,d-1\},
\qquad
S(0)=S(1)=Y(0)=0,
\qquad
Y(1)\mid X\sim\operatorname{Bernoulli}\!\left(\frac q{1+q}\right).
\]
Draw \(A\sim\operatorname{Bernoulli}(1/2)\) independently of the baseline and potential variables, set
\[
S=S(A),\qquad Y=Y(A),
\]
and draw the arrival indicator conditionally independently of \(Y(1)\), with
\[
P_\star(R=1\mid X,A,Y(1))
=
\begin{cases}
1,&A=0,\\
(1+q)/2,&A=1.
\end{cases}
\]
These distributions are well defined because
\[
0<\frac q{1+q}\le\frac12,
\qquad
q\le\frac{1+q}{2}\le1.
\]
Evaluating the average treatment effect of
\cref{obj:def:ate-functional} yields
\[
\tau(P_\star)
=E_{P_\star}[Y(1)-Y(0)]
=\frac q{1+q}
=\frac12-h_{\mathrm{sep}}.
\]
% lean: CausalSmith.Stat.MarRareqLogfrontier.largeAlphabet_center_separation
% lean: CausalSmith.Stat.MarRareqLogfrontier.largeAlphabetComparisonLaw_ate

\item For each
\[
\xi=(\xi_x)_{x=0}^{d-1}\in\{0,1\}^{\{0,\ldots,d-1\}},
\]
construct a law \(P_\xi\) with the same uniform baseline and independent fair treatment assignment, and with
\[
S(0)=S(1)=Y(0)=0,\qquad Y(1)=\xi_X,
\qquad S=S(A),\qquad Y=Y(A).
\]
Its arrival draw satisfies, for \(x\in\{0,\ldots,d-1\}\) and \(a\in\{0,1\}\),
\[
P_\xi(R=1\mid X=x,A=a)=1-a\delta\xi_x.
\]
Thus every arrival probability is either \(q\) or \(1\), and
\[
\tau(P_\xi)
=\frac1d\sum_{x=0}^{d-1}\xi_x.
\]
Put the product fair-coin prior on these vectors:
\[
\nu_{\mathrm{cell}}
:=\bigotimes_{x=0}^{d-1}\operatorname{Bernoulli}(1/2),
\qquad
\nu_{\mathrm{cell}}(\{\xi\})=2^{-d}.
\]
Define the sample laws and their mixture by
\[
Q_\xi:=\mathcal L_{P_\xi}(\mathbf O_n),
\qquad
Q_\star:=\mathcal L_{P_\star}(\mathbf O_n),
\qquad
Q_{\mathrm{mix}}
:=\int Q_\xi\,\nu_{\mathrm{cell}}(d\xi)
=\sum_\xi2^{-d}Q_\xi.
\]
For each fixed full-data law, the sample consists of \(n\) independent observed records as in \cref{obj:ass:iid-sampling}. The mixture draws \(\xi\) once for the entire sample.

Define the test event, the exceptional prior vectors, and the two risks by
\[
E_{\mathrm{test}}
:=
\left\{(\mathbf o,w):
w\le\frac12-\frac{h_{\mathrm{sep}}}{2}\right\},
\qquad
w\in\mathbb R,
\]
\[
F_{\mathrm{bad}}
:=
\left\{\xi:
\left|\tau(P_\xi)-\frac12\right|>
\frac{h_{\mathrm{sep}}}{4}\right\},
\]
\[
\mathcal R_{\mathrm{prior}}(\mathsf T)
:=\int\mathcal R_{\mathsf T}(P_\xi)\,\nu_{\mathrm{cell}}(d\xi),
\qquad
\mathcal R_\star(\mathsf T):=\mathcal R_{\mathsf T}(P_\star).
\]
The notation \(Q\ltimes\mathsf T\) denotes the joint sample--action law used in
\cref{obj:lem:randomized-kernel-tv-comparison}.

If \(\xi\notin F_{\mathrm{bad}}\) and \((\mathbf o,w)\in E_{\mathrm{test}}\), then
\[
\tau(P_\xi)-w
\ge
\left(\frac12-\frac{h_{\mathrm{sep}}}{4}\right)
-\left(\frac12-\frac{h_{\mathrm{sep}}}{2}\right)
=\frac{h_{\mathrm{sep}}}{4}.
\]
Consequently, for every \(\xi\),
\[
\frac{h_{\mathrm{sep}}^2}{16}
(Q_\xi\ltimes\mathsf T)(E_{\mathrm{test}})
\le
\mathcal R_{\mathsf T}(P_\xi)
+\frac{h_{\mathrm{sep}}^2}{16}
\mathbf1_{F_{\mathrm{bad}}}(\xi).
\]
For exceptional vectors this follows from the probability bound by one and nonnegativity of risk; for the other vectors it follows by integrating the displayed pointwise loss bound. Since the same kernel is used for every \(\xi\), finite averaging gives
\[
(Q_{\mathrm{mix}}\ltimes\mathsf T)(E_{\mathrm{test}})
=
\int
(Q_\xi\ltimes\mathsf T)(E_{\mathrm{test}})
\,\nu_{\mathrm{cell}}(d\xi).
\]
Averaging the member-wise risk inequality therefore yields
\[
\frac{h_{\mathrm{sep}}^2}{16}
(Q_{\mathrm{mix}}\ltimes\mathsf T)(E_{\mathrm{test}})
\le
\mathcal R_{\mathrm{prior}}(\mathsf T)
+\frac{h_{\mathrm{sep}}^2}{16}\nu_{\mathrm{cell}}(F_{\mathrm{bad}}).
\]

Under \(P_\star\), an action in \(E_{\mathrm{test}}^c\) satisfies
\[
w-\tau(P_\star)
>
\left(\frac12-\frac{h_{\mathrm{sep}}}{2}\right)
-\left(\frac12-h_{\mathrm{sep}}\right)
=\frac{h_{\mathrm{sep}}}{2}.
\]
In particular, its squared loss is at least \(h_{\mathrm{sep}}^2/16\), so
\[
\frac{h_{\mathrm{sep}}^2}{16}
(Q_\star\ltimes\mathsf T)(E_{\mathrm{test}}^c)
\le\mathcal R_\star(\mathsf T).
\]
% lean: CausalSmith.Stat.MarRareqLogfrontier.largeAlphabetLaw_ate
% lean: CausalSmith.Stat.MarRareqLogfrontier.largeAlphabet_prior_test_risk
% lean: CausalSmith.Stat.MarRareqLogfrontier.largeAlphabet_comparison_test_risk

\item We next compare the two sample laws. For a single treated visit to any baseline cell \(x\), averaging the fair coordinate \(\xi_x\) gives the following observed-mark probabilities:
\[
\begin{array}{c|c|c}
(R,RY)
&
\displaystyle\int P_\xi((R,RY)\mid X=x,A=1)\,
\nu_{\mathrm{cell}}(d\xi)
&
P_\star((R,RY)\mid X=x,A=1)
\\ \hline
(1,1)&q/2&
\displaystyle\frac{1+q}{2}\frac q{1+q}=q/2
\\
(1,0)&1/2&
\displaystyle\frac{1+q}{2}\left(1-\frac q{1+q}\right)=1/2
\\
(0,0)&\delta/2&
\displaystyle1-\frac{1+q}{2}=\delta/2
\\
(0,1)&0&0
\end{array}
\]
In the prior column, the first probability comes from \(\xi_x=1\) and arrival, the second from \(\xi_x=0\), and the third from \(\xi_x=1\) and nonarrival. A control visit has mark \((R,RY)=(1,0)\) with probability one under every law. All surrogates equal zero.

Define the event of distinct treated baseline cells by
\[
E_{\mathrm{distinct}}
:=
\bigcap_{1\le i<j\le n}
\{A_i\ne1\ \text{or}\ A_j\ne1\ \text{or}\ X_i\ne X_j\}.
\]
For any sample atom \(\mathbf o=(o_1,\ldots,o_n)\) in this event, each treated-record probability depends on a different prior coordinate, while each control-record probability is constant in \(\xi\). Independence of the prior coordinates and the one-visit calculation therefore give
\[
\begin{aligned}
Q_{\mathrm{mix}}(\{\mathbf o\})
&=
\int
\prod_{i=1}^n
\mathcal L_{P_\xi}(O_1)(\{o_i\})\,
\nu_{\mathrm{cell}}(d\xi)
\\
&=
\prod_{i=1}^n
\int
\mathcal L_{P_\xi}(O_1)(\{o_i\})\,
\nu_{\mathrm{cell}}(d\xi)
\\
&=
\prod_{i=1}^n
\mathcal L_{P_\star}(O_1)(\{o_i\})
=
Q_\star(\{\mathbf o\}).
\end{aligned}
\]
The equality also covers atoms with impossible surrogate or arrived-outcome marks, whose probabilities are zero under both constructions. % lean: CausalSmith.Stat.MarRareqLogfrontier.largeAlphabetSampleMixture_matching

\item Under each \(P_\xi\) and under \(P_\star\), the probability of any treated baseline cell is
\[
P_\xi(A=1,X=x)=P_\star(A=1,X=x)=\frac1{2d},
\qquad x\in\{0,\ldots,d-1\}.
\]
Thus independence of distinct observations gives, for \(1\le i<j\le n\),
\[
\begin{aligned}
Q_\xi\{A_i=A_j=1,\ X_i=X_j\}
&\le
\sum_{x=0}^{d-1}\left(\frac1{2d}\right)^2
=\frac1{4d},
\\
Q_\star\{A_i=A_j=1,\ X_i=X_j\}
&\le\frac1{4d}.
\end{aligned}
\]
Taking a union over the \(\binom n2\) pairs yields
\[
Q_\xi(E_{\mathrm{distinct}}^c),
\ Q_\star(E_{\mathrm{distinct}}^c)
\le
\binom n2\frac1{4d}
\le\frac{n^2}{8d}
\le\frac1{8192}.
\]
When \(n=1\), the pair set is empty and these collision probabilities are zero. Averaging the first bound over the prior gives
\[
Q_{\mathrm{mix}}(E_{\mathrm{distinct}}^c)
=
\int Q_\xi(E_{\mathrm{distinct}}^c)\,
\nu_{\mathrm{cell}}(d\xi)
\le\frac1{8192}.
\]

On the finite sample space, agreement of the atom masses on
\(E_{\mathrm{distinct}}\) implies
\[
\begin{aligned}
\operatorname{TV}(Q_{\mathrm{mix}},Q_\star)
&\le
\frac12\sum_{\mathbf o}
\left|Q_{\mathrm{mix}}(\{\mathbf o\})
-Q_\star(\{\mathbf o\})\right|
\\
&\le
\frac12\left\{
Q_{\mathrm{mix}}(E_{\mathrm{distinct}}^c)
+Q_\star(E_{\mathrm{distinct}}^c)
\right\}
\le\frac1{8192}.
\end{aligned}
\]
Indeed, the summands vanish on \(E_{\mathrm{distinct}}\), and elsewhere each absolute difference is bounded by the sum of the two atom masses. Applying
\cref{obj:lem:randomized-kernel-tv-comparison} to the common kernel and the measurable event \(E_{\mathrm{test}}\) now gives
\[
1-\frac1{8192}
\le
1-\operatorname{TV}(Q_{\mathrm{mix}},Q_\star)
\le
(Q_{\mathrm{mix}}\ltimes\mathsf T)(E_{\mathrm{test}})
+
(Q_\star\ltimes\mathsf T)(E_{\mathrm{test}}^c).
\]
% lean: CausalSmith.Stat.MarRareqLogfrontier.largeAlphabet_actual_collision_testing

\item Under \(\nu_{\mathrm{cell}}\), the independent fair coordinates have
\[
E_{\nu_{\mathrm{cell}}}[\xi_x]=\frac12,
\qquad
\operatorname{Var}_{\nu_{\mathrm{cell}}}(\xi_x)=\frac14.
\]
Consequently,
\[
E_{\nu_{\mathrm{cell}}}\!\left[\sum_{x=0}^{d-1}\xi_x\right]
=\frac d2,
\qquad
\operatorname{Var}_{\nu_{\mathrm{cell}}}
\!\left(\sum_{x=0}^{d-1}\xi_x\right)
=\frac d4.
\]
The regime \(n^{-1}\le\delta^2\), together with \(\delta\ge0\), implies
\[
\delta>0,\qquad h_{\mathrm{sep}}\ge\frac{\delta}{4}>0,
\qquad
1\le n\delta^2,\qquad
\delta^2\le16h_{\mathrm{sep}}^2.
\]
These inequalities and the alphabet-size condition give the concentration budget
\[
d h_{\mathrm{sep}}^2
\ge1024n^2h_{\mathrm{sep}}^2
\ge64n^2\delta^2
\ge64n
\ge64.
\]
Since \(\tau(P_\xi)=d^{-1}\sum_x\xi_x\), Chebyshev's inequality applied to the sum of the prior coordinates yields
\[
\begin{aligned}
\nu_{\mathrm{cell}}(F_{\mathrm{bad}})
&=
\nu_{\mathrm{cell}}
\left\{
\left|\sum_{x=0}^{d-1}\xi_x-\frac d2\right|
>\frac{d h_{\mathrm{sep}}}{4}
\right\}
\\
&\le
\frac{d/4}{(d h_{\mathrm{sep}}/4)^2}
=
\frac4{d h_{\mathrm{sep}}^2}
\le\frac1{16}.
\end{aligned}
\]
% lean: CausalSmith.Stat.MarRareqLogfrontier.largeAlphabet_center_prior_concentration

\item Multiplying the testing and concentration bounds by the nonnegative factor \(h_{\mathrm{sep}}^2/16\) gives
\[
\frac{h_{\mathrm{sep}}^2}{16}
\left(1-\frac1{8192}\right)
\le
\frac{h_{\mathrm{sep}}^2}{16}
\left\{
(Q_{\mathrm{mix}}\ltimes\mathsf T)(E_{\mathrm{test}})
+
(Q_\star\ltimes\mathsf T)(E_{\mathrm{test}}^c)
\right\},
\]
\[
\frac{h_{\mathrm{sep}}^2}{16}
\nu_{\mathrm{cell}}(F_{\mathrm{bad}})
\le\frac{h_{\mathrm{sep}}^2}{256}.
\]
Combining these with the two test-risk inequalities established above, we obtain
\[
\begin{aligned}
\mathcal R_{\mathrm{prior}}(\mathsf T)
+\mathcal R_\star(\mathsf T)
&\ge
\frac{h_{\mathrm{sep}}^2}{16}
\left\{
1-\frac1{8192}-\frac1{16}
\right\}
\\
&\ge\frac{h_{\mathrm{sep}}^2}{32},
\end{aligned}
\]
where the last inequality uses
\(1-1/8192-1/16\ge1/2\). % lean: CausalSmith.Stat.MarRareqLogfrontier.largeAlphabet_deficit_minimax_lower

\item To pass to worst-case risk, define
\[
\mathcal R_{\mathrm{sup}}(\mathsf T)
:=
\sup_{P\in\mathcal M^{+}_{n,d,q}}
\mathcal R_{\mathsf T}(P).
\]
Binary potential outcomes give \(\tau(P)\in[-1,1]\). Since the kernel takes values in \([-1,1]\), its loss and risk satisfy
\[
0\le(w-\tau(P))^2\le4
\quad(w\in[-1,1]),
\qquad
0\le\mathcal R_{\mathsf T}(P)\le4.
\]
Thus the member-law risks are bounded above.

Every constructed law belongs to the model class in
\cref{obj:def:unrestricted-arrival-model-class}. The independent fair treatment draw verifies
\cref{obj:ass:randomized-independence,obj:ass:balanced-randomization}, and the defining identities for the realized variables verify
\cref{obj:ass:surrogate-consistency,obj:ass:outcome-consistency}.
Within an occupied cell, the outcome under \(P_\xi\) is deterministic, while under \(P_\star\) its arrival draw is independent of the outcome; these properties verify
\cref{obj:ass:arrival-mar}. For every \(a\in\{0,1\}\) and \(x\in\{0,\ldots,d-1\}\), the cell probabilities of
\cref{obj:synth_16} are
\[
p_{ax0}(P_\xi)=p_{ax0}(P_\star)=\frac1{2d},
\qquad
p_{ax1}(P_\xi)=p_{ax1}(P_\star)=0.
\]
On the occupied cells, their arrival probabilities are
\[
\rho_{ax0}(P_\xi)=1-a\delta\xi_x\ge q,
\qquad
\rho_{ax0}(P_\star)
=
\begin{cases}
1,&a=0,\\
(1+q)/2,&a=1,
\end{cases}
\ge q.
\]
This verifies \cref{obj:ass:occupied-cell-arrival}, and hence
\[
P_\xi,P_\star\in\mathcal M^{+}_{n,d,q}.
\]

Each member-law risk is therefore at most
\(\mathcal R_{\mathrm{sup}}(\mathsf T)\). Since
\(\nu_{\mathrm{cell}}\) is a probability law,
\[
\mathcal R_{\mathrm{prior}}(\mathsf T)
\le
\int\mathcal R_{\mathrm{sup}}(\mathsf T)\,
\nu_{\mathrm{cell}}(d\xi)
=
\mathcal R_{\mathrm{sup}}(\mathsf T),
\qquad
\mathcal R_\star(\mathsf T)
\le\mathcal R_{\mathrm{sup}}(\mathsf T).
\]
It follows that
\[
2\mathcal R_{\mathrm{sup}}(\mathsf T)
\ge
\mathcal R_{\mathrm{prior}}(\mathsf T)+\mathcal R_\star(\mathsf T)
\ge\frac{h_{\mathrm{sep}}^2}{32}.
\]
Using \(h_{\mathrm{sep}}\ge\delta/4\ge0\), we conclude
\[
\mathcal R_{\mathrm{sup}}(\mathsf T)
\ge\frac{h_{\mathrm{sep}}^2}{64}
\ge\frac{\delta^2}{1024}.
\]
This holds for every kernel in the stated decision class. Taking its infimum, as in
\cref{obj:def:unrestricted-minimax-risk}, proves
\[
\frac{\delta^2}{1024}\le\mathfrak R^{+}_{n,d,q}
\]
in the remaining regime. % lean: CausalSmith.Stat.MarRareqLogfrontier.largeAlphabet_deficit_minimax_lower

\item The two cases establish the deficit lower bound throughout the stated parameter range. Combining it with the initial bounds gives
\[
\max\left\{\frac1{256n},\frac{\delta^2}{1024}\right\}
\le\mathfrak R^{+}_{n,d,q}
\le\frac4n+\delta^2.
\]
Adding the two lower inequalities and dividing by two yields
\[
\mathfrak R^{+}_{n,d,q}
\ge
\frac1{512n}+\frac{\delta^2}{2048}
\ge
\frac1{2048}\left\{\frac1n+\delta^2\right\}.
\]
Finally, \(\delta^2\ge0\) gives
\[
\mathfrak R^{+}_{n,d,q}
\le\frac4n+\delta^2
\le4\left\{\frac1n+\delta^2\right\}.
\]
Substituting \(\delta=1-q\) proves the asserted envelope. % lean: CausalSmith.Stat.MarRareqLogfrontier.unrestricted_large_alphabet_deficit_lower
\end{enumerate}
\end{proof}
